% figures/ch6/lo_feynman_diagrams.tex
% 2x2 grid of leading-order Feynman diagrams for V+jets (V = Z or W) production.
% Included via \input inside a figure environment in chapter6.tex.

% ---------------------------------------------------------------
% Row 1: quark--antiquark annihilation channel  qq' -> Vg
% ---------------------------------------------------------------
\begin{subfigure}[b]{0.46\textwidth}
  \centering
  \begin{tikzpicture}
    \begin{feynman}
      % Incoming legs enter horizontally from the left
      \vertex (iq)  at (0.0,  1.0) {\(q\)};
      \vertex (iqb) at (0.0, -1.0) {\(\bar{q}'\)};
      % Interaction vertices: v1 on the upper rail, v2 at centre-right
      \vertex (v1)  at (2.0,  1.0);
      \vertex (v2)  at (2.0, -1.0);
      % Outgoing: V exits upper-right at 45 deg, gluon exits horizontally
      \vertex (fV)  at (4.0,  1.0) {\(V\)};
      \vertex (fg)  at (4.0, -1.0) {\(g\)};
      \diagram* {
        (iq)  -- [fermion]      (v1),
        (iqb) -- [anti fermion] (v2),
        (v1)  -- [boson]        (fV),
        (v1)  -- [fermion]      (v2),
        (v2)  -- [gluon]        (fg),
      };
    \end{feynman}
  \end{tikzpicture}
  \caption*{(a)}
\end{subfigure}
\hfill
\begin{subfigure}[b]{0.46\textwidth}
  \centering
  \begin{tikzpicture}
    \begin{feynman}
      % Mirror of (a): V emitted from antiquark rail instead
      \vertex (iq)  at (0.0,  1.0) {\(q\)};
      \vertex (iqb) at (0.0, -1.0) {\(\bar{q}'\)};
      \vertex (v1)  at (2.0,  1.0);
      \vertex (v2)  at (2.0, -1.0);
      \vertex (fV)  at (4.0, -1.0) {\(V\)};
      \vertex (fg)  at (4.0,  1.0) {\(g\)};
      \diagram* {
        (iq)  -- [fermion]      (v1),
        (iqb) -- [anti fermion] (v2),
        (v2)  -- [boson]        (fV),
        (v1)  -- [fermion]      (v2),
        (v1)  -- [gluon]        (fg),
      };
    \end{feynman}
  \end{tikzpicture}
  \caption*{(b)}
\end{subfigure}

\vspace{2em}

% ---------------------------------------------------------------
% Row 2: quark--gluon Compton channel  qg -> Vq
% ---------------------------------------------------------------
\begin{subfigure}[b]{0.46\textwidth}
  \centering
  \begin{tikzpicture}
    \begin{feynman}
      % t-channel variant: g absorbed first, then V emitted
      \vertex (iq)  at (0.0,  1.0) {\(q\)};
      \vertex (ig)  at (0.0, -1.0) {\(g\)};
      \vertex (v1)  at (1.33,  0.0);
      \vertex (v2)  at (2.66,  0.0);
      \vertex (fV)  at (4.0,  1.0) {\(V\)};
      \vertex (fq)  at (4.0, -1.0) {\(q\)};
      \diagram* {
        (iq)  -- [fermion] (v1),
        (ig)  -- [gluon]   (v1),
        (v1)  -- [fermion] (v2),
        (v2)  -- [boson]   (fV),
        (v2)  -- [fermion] (fq),
      };
    \end{feynman}
  \end{tikzpicture}
  \caption*{(c)}
\end{subfigure}
\hfill
\begin{subfigure}[b]{0.46\textwidth}
  \centering
  \begin{tikzpicture}
    \begin{feynman}
      % t-channel variant: g absorbed first, then V emitted
      \vertex (iq)  at (0.0,  1.0) {\(q\)};
      \vertex (ig)  at (0.0, -1.0) {\(g\)};
      \vertex (v1)  at (2.0,  1.0);
      \vertex (v2)  at (2.0,  -1.0);
      \vertex (fV)  at (4.0,  1.0) {\(V\)};
      \vertex (fq)  at (4.0, -1.0) {\(q\)};
      \diagram* {
        (iq)  -- [fermion] (v1),
        (ig)  -- [gluon]   (v2),
        (v1)  -- [fermion] (v2),
        (v1)  -- [boson]   (fV),
        (v2)  -- [fermion] (fq),
      };
    \end{feynman}
  \end{tikzpicture}
  \caption*{(d)}
\end{subfigure}